% sections/ablation.tex -- Section 8: ablations A1, A4, A5 (W25).
% Source: experiments/ablation/REPORT.md (W18); fig_ablation_a1.pdf by W26 (figs/FIGURES.md §4).
\section{What the Proof Should Bind: Ablations}\label{sec:ablation}

Binding more of the request is an old idea---the 2016 OAuth signed-request
draft covered query, headers and body~\cite{p2:richer2016signedhttp}---and
RFC~9421 makes coverage configurable~\cite{rfc9421}. What has been missing is
evidence of what each claim actually buys in a layered deployment. Ablation
\Abl{1} provides it, and directly informs the current proposal to make \htm and
\htu optional~\cite{allam2026thread}.

\paragraph{Design (\Abl{1}).} A copy of the reference verifier implements
\AblOneVariants{} claim sets: \emph{full} (RFC~9449); \emph{$-$\htm},
\emph{$-$\htu}, \emph{$-$\ath} (claim neither required nor checked);
\emph{+query} (\htu includes the RFC~3986-normalised query; an absent query is
bound too); \emph{+digest} (a proof claim must equal the request's RFC~9530
\hdr{Content-Digest}~\cite{w2:rfc9530}, and a body-holding hop between NGINX and
the executor rejects any body that does not match it); \emph{digest-nogate}
(the same without that hop); and \emph{+query+digest}. The two halves of
+digest are both needed: NGINX \texttt{auth\_request} sends the subrequest
without the body, so an edge verifier never sees the content. We ran
\AblOneTopos{} topologies (NGINX default, four frameworks with v2 routes,
override off/on, $\times$ \AblOneVariants{} variants): \AblOneRequests{}
mutation requests from \AblOneCases{} cases per topology, plus
\AblOneSwapRequests{} \emph{swap} requests that replay a proof for operation
$p$ on the canonical request of operation $q$ or with another method---a case
the mutation catalogue never generates. Hits were confirmed 3/3.

\begin{figure*}[t]
\centering
\figpdf{\textwidth}{4.3cm}{fig_ablation_a1.pdf}
\Description{Left: stacked bars of residual hits on the mutation set for eight claim variants; full, no-htm, no-ath and digest without gate at 192, no-htu at 432, plus-query 147, plus-digest 129, plus-query-plus-digest 84 with a dashed line at the 60 G2 hits that no variant removes. Right: cross-operation and method-swap replay hits, zero for all variants except no-htm at 192 and no-htu at 720.}
\caption{\Abl{1} binding-claim ablation. (a)~Residual hits on the mutation set:
\htu does nearly all the work; binding the query and the body (with a
body-holding hop) removes the dispatch hits, but \AblOneHitsQueryDigest{}
remain: header-carried \GA and the \AblOneGTwoSurvive{} \GB hits that no claim
set removes (dashed). (b)~Cross-operation and method-swap replay: \htm and \htu
each prevent replay that the mutation corpus cannot see.}
\label{fig:a1}
\end{figure*}

\paragraph{Results.} Figure~\ref{fig:a1} summarises them.
(1)~\emph{\htu does nearly all the work against mutation:} every verifier
denial on the mutation set is an \htu mismatch; dropping it raises the hits
from \AblOneHitsFull{} to \AblOneHitsNoHtu{} and resolution from
\AblOneResFull{} to \AblOneResNoHtu{}, and admits \AblOneSwapNoHtu{}
same-method cross-operation replays (resolution \AblOneSwapResNoHtu{}).
(2)~\emph{\htm is invisible to path mutation but necessary:} $-$\htm gives the
identical hit set on the mutation corpus, yet admits \AblOneSwapNoHtm{}
same-path method swaps (resolution \AblOneSwapResNoHtm{}). A corpus without
method swaps would wrongly conclude that \htm is redundant.
(3)~\emph{\ath buys nothing under \Aleak}, as expected: the attacker replays the
proof with its own token.
(4)~\emph{+query removes exactly the query class} (\AblOneQueryRemoved{} hits,
including \texttt{?\_method}) and adds none, but inherits the fragment
differential: \texttt{?op=freeze\#frag} still passes (\AblOneFragQuery{}
hits), because the verifier puts \texttt{\#frag} in the fragment while three
servers keep it in the query value.
(5)~\emph{+digest removes the body classes only with a body-holding hop}
(\AblOneDigestRemoved{} hits; the hop rejected \AblOneGateDeny{}
body-mismatched requests); \emph{digest-nogate} gives the same hits as
\emph{full}, because the attacker keeps the leaked header and swaps the body.
(6)~\emph{Binding the query and the body does not close everything.}
+query+digest leaves \AblOneHitsQueryDigest{} hits: header-carried \GA
(\hdr{X-HTTP-Method-Override}, the \hdr{Content-Type} relabel, \texttt{\#} in
the query) and all \AblOneGTwoSurvive{} \GB path-differential hits, which no
claim set removes. Covering more of the request does not fix interpretation
differentials; only agreement between the verifier's and the router's view
does---the argument for the executor-view contract and D-B
(\S\ref{sec:defence}).
\emph{Legitimate traffic.} Across \AblOneLegitN{} legitimate requests,
query binding caused no false rejection behind NGINX, which forwards the raw
target; behind HAProxy's recommended query normalisers, exact matching
rejected \AblOneLegitFRHapExact{}, RFC~3986 matching
\AblOneLegitFRHapRfc{}, and sorted matching \AblOneLegitFRHapSorted{}. Body
binding behind HAProxy failed for every request because its adapter does not
forward \hdr{Content-Digest}: the risk is integration, not transformation.

\paragraph{Primitive families (\Abl{4}) and profiles (\Abl{5}).}
Leaving one primitive family out shows that \AblFourResult{}. The \GB evidence
is therefore fragile to catalogue coverage: dropping the fragment family alone
removes B1. Across profiles, \AblFiveResult{} (\AblFiveGTwoTopos{} new \GB and
\AblFiveEFourTopos{} new E4 topologies). The strict deny rule lists
\texttt{\%23} but not a raw \texttt{\#}, and strict profiles are
harness-defined, so they are a ceiling rather than a vendor default.
